% Options for packages loaded elsewhere
% Options for packages loaded elsewhere
\PassOptionsToPackage{unicode}{hyperref}
\PassOptionsToPackage{hyphens}{url}
\PassOptionsToPackage{dvipsnames,svgnames,x11names}{xcolor}
%
\documentclass[
]{report}
\usepackage{xcolor}
\usepackage[lmargin=30mm,rmargin=30mm]{geometry}
\usepackage{amsmath,amssymb}
\setcounter{secnumdepth}{-\maxdimen} % remove section numbering
\usepackage{iftex}
\ifPDFTeX
  \usepackage[T1]{fontenc}
  \usepackage[utf8]{inputenc}
  \usepackage{textcomp} % provide euro and other symbols
\else % if luatex or xetex
  \usepackage{unicode-math} % this also loads fontspec
  \defaultfontfeatures{Scale=MatchLowercase}
  \defaultfontfeatures[\rmfamily]{Ligatures=TeX,Scale=1}
\fi
\usepackage{lmodern}
\ifPDFTeX\else
  % xetex/luatex font selection
\fi
% Use upquote if available, for straight quotes in verbatim environments
\IfFileExists{upquote.sty}{\usepackage{upquote}}{}
\IfFileExists{microtype.sty}{% use microtype if available
  \usepackage[]{microtype}
  \UseMicrotypeSet[protrusion]{basicmath} % disable protrusion for tt fonts
}{}
\makeatletter
\@ifundefined{KOMAClassName}{% if non-KOMA class
  \IfFileExists{parskip.sty}{%
    \usepackage{parskip}
  }{% else
    \setlength{\parindent}{0pt}
    \setlength{\parskip}{6pt plus 2pt minus 1pt}}
}{% if KOMA class
  \KOMAoptions{parskip=half}}
\makeatother
% Make \paragraph and \subparagraph free-standing
\makeatletter
\ifx\paragraph\undefined\else
  \let\oldparagraph\paragraph
  \renewcommand{\paragraph}{
    \@ifstar
      \xxxParagraphStar
      \xxxParagraphNoStar
  }
  \newcommand{\xxxParagraphStar}[1]{\oldparagraph*{#1}\mbox{}}
  \newcommand{\xxxParagraphNoStar}[1]{\oldparagraph{#1}\mbox{}}
\fi
\ifx\subparagraph\undefined\else
  \let\oldsubparagraph\subparagraph
  \renewcommand{\subparagraph}{
    \@ifstar
      \xxxSubParagraphStar
      \xxxSubParagraphNoStar
  }
  \newcommand{\xxxSubParagraphStar}[1]{\oldsubparagraph*{#1}\mbox{}}
  \newcommand{\xxxSubParagraphNoStar}[1]{\oldsubparagraph{#1}\mbox{}}
\fi
\makeatother

\usepackage{color}
\usepackage{fancyvrb}
\newcommand{\VerbBar}{|}
\newcommand{\VERB}{\Verb[commandchars=\\\{\}]}
\DefineVerbatimEnvironment{Highlighting}{Verbatim}{commandchars=\\\{\}}
% Add ',fontsize=\small' for more characters per line
\usepackage{framed}
\definecolor{shadecolor}{RGB}{241,243,245}
\newenvironment{Shaded}{\begin{snugshade}}{\end{snugshade}}
\newcommand{\AlertTok}[1]{\textcolor[rgb]{0.68,0.00,0.00}{#1}}
\newcommand{\AnnotationTok}[1]{\textcolor[rgb]{0.37,0.37,0.37}{#1}}
\newcommand{\AttributeTok}[1]{\textcolor[rgb]{0.40,0.45,0.13}{#1}}
\newcommand{\BaseNTok}[1]{\textcolor[rgb]{0.68,0.00,0.00}{#1}}
\newcommand{\BuiltInTok}[1]{\textcolor[rgb]{0.00,0.23,0.31}{#1}}
\newcommand{\CharTok}[1]{\textcolor[rgb]{0.13,0.47,0.30}{#1}}
\newcommand{\CommentTok}[1]{\textcolor[rgb]{0.37,0.37,0.37}{#1}}
\newcommand{\CommentVarTok}[1]{\textcolor[rgb]{0.37,0.37,0.37}{\textit{#1}}}
\newcommand{\ConstantTok}[1]{\textcolor[rgb]{0.56,0.35,0.01}{#1}}
\newcommand{\ControlFlowTok}[1]{\textcolor[rgb]{0.00,0.23,0.31}{\textbf{#1}}}
\newcommand{\DataTypeTok}[1]{\textcolor[rgb]{0.68,0.00,0.00}{#1}}
\newcommand{\DecValTok}[1]{\textcolor[rgb]{0.68,0.00,0.00}{#1}}
\newcommand{\DocumentationTok}[1]{\textcolor[rgb]{0.37,0.37,0.37}{\textit{#1}}}
\newcommand{\ErrorTok}[1]{\textcolor[rgb]{0.68,0.00,0.00}{#1}}
\newcommand{\ExtensionTok}[1]{\textcolor[rgb]{0.00,0.23,0.31}{#1}}
\newcommand{\FloatTok}[1]{\textcolor[rgb]{0.68,0.00,0.00}{#1}}
\newcommand{\FunctionTok}[1]{\textcolor[rgb]{0.28,0.35,0.67}{#1}}
\newcommand{\ImportTok}[1]{\textcolor[rgb]{0.00,0.46,0.62}{#1}}
\newcommand{\InformationTok}[1]{\textcolor[rgb]{0.37,0.37,0.37}{#1}}
\newcommand{\KeywordTok}[1]{\textcolor[rgb]{0.00,0.23,0.31}{\textbf{#1}}}
\newcommand{\NormalTok}[1]{\textcolor[rgb]{0.00,0.23,0.31}{#1}}
\newcommand{\OperatorTok}[1]{\textcolor[rgb]{0.37,0.37,0.37}{#1}}
\newcommand{\OtherTok}[1]{\textcolor[rgb]{0.00,0.23,0.31}{#1}}
\newcommand{\PreprocessorTok}[1]{\textcolor[rgb]{0.68,0.00,0.00}{#1}}
\newcommand{\RegionMarkerTok}[1]{\textcolor[rgb]{0.00,0.23,0.31}{#1}}
\newcommand{\SpecialCharTok}[1]{\textcolor[rgb]{0.37,0.37,0.37}{#1}}
\newcommand{\SpecialStringTok}[1]{\textcolor[rgb]{0.13,0.47,0.30}{#1}}
\newcommand{\StringTok}[1]{\textcolor[rgb]{0.13,0.47,0.30}{#1}}
\newcommand{\VariableTok}[1]{\textcolor[rgb]{0.07,0.07,0.07}{#1}}
\newcommand{\VerbatimStringTok}[1]{\textcolor[rgb]{0.13,0.47,0.30}{#1}}
\newcommand{\WarningTok}[1]{\textcolor[rgb]{0.37,0.37,0.37}{\textit{#1}}}

\usepackage{longtable,booktabs,array}
\newcounter{none} % for unnumbered tables
\usepackage{calc} % for calculating minipage widths
% Correct order of tables after \paragraph or \subparagraph
\usepackage{etoolbox}
\makeatletter
\patchcmd\longtable{\par}{\if@noskipsec\mbox{}\fi\par}{}{}
\makeatother
% Allow footnotes in longtable head/foot
\IfFileExists{footnotehyper.sty}{\usepackage{footnotehyper}}{\usepackage{footnote}}
\makesavenoteenv{longtable}
\usepackage{graphicx}
\makeatletter
\newsavebox\pandoc@box
\newcommand*\pandocbounded[1]{% scales image to fit in text height/width
  \sbox\pandoc@box{#1}%
  \Gscale@div\@tempa{\textheight}{\dimexpr\ht\pandoc@box+\dp\pandoc@box\relax}%
  \Gscale@div\@tempb{\linewidth}{\wd\pandoc@box}%
  \ifdim\@tempb\p@<\@tempa\p@\let\@tempa\@tempb\fi% select the smaller of both
  \ifdim\@tempa\p@<\p@\scalebox{\@tempa}{\usebox\pandoc@box}%
  \else\usebox{\pandoc@box}%
  \fi%
}
% Set default figure placement to htbp
\def\fps@figure{htbp}
\makeatother

\ifLuaTeX
  \usepackage{luacolor}
  \usepackage[soul]{lua-ul}
\else
  \usepackage{soul}
\fi




\setlength{\emergencystretch}{3em} % prevent overfull lines

\providecommand{\tightlist}{%
  \setlength{\itemsep}{0pt}\setlength{\parskip}{0pt}}



 


\makeatletter
\@ifpackageloaded{tcolorbox}{}{\usepackage[skins,breakable]{tcolorbox}}
\@ifpackageloaded{fontawesome5}{}{\usepackage{fontawesome5}}
\definecolor{quarto-callout-color}{HTML}{909090}
\definecolor{quarto-callout-note-color}{HTML}{0758E5}
\definecolor{quarto-callout-important-color}{HTML}{CC1914}
\definecolor{quarto-callout-warning-color}{HTML}{EB9113}
\definecolor{quarto-callout-tip-color}{HTML}{00A047}
\definecolor{quarto-callout-caution-color}{HTML}{FC5300}
\definecolor{quarto-callout-color-frame}{HTML}{acacac}
\definecolor{quarto-callout-note-color-frame}{HTML}{4582ec}
\definecolor{quarto-callout-important-color-frame}{HTML}{d9534f}
\definecolor{quarto-callout-warning-color-frame}{HTML}{f0ad4e}
\definecolor{quarto-callout-tip-color-frame}{HTML}{02b875}
\definecolor{quarto-callout-caution-color-frame}{HTML}{fd7e14}
\makeatother
\makeatletter
\@ifpackageloaded{caption}{}{\usepackage{caption}}
\AtBeginDocument{%
\ifdefined\contentsname
  \renewcommand*\contentsname{Table of contents}
\else
  \newcommand\contentsname{Table of contents}
\fi
\ifdefined\listfigurename
  \renewcommand*\listfigurename{List of Figures}
\else
  \newcommand\listfigurename{List of Figures}
\fi
\ifdefined\listtablename
  \renewcommand*\listtablename{List of Tables}
\else
  \newcommand\listtablename{List of Tables}
\fi
\ifdefined\figurename
  \renewcommand*\figurename{Figure}
\else
  \newcommand\figurename{Figure}
\fi
\ifdefined\tablename
  \renewcommand*\tablename{Table}
\else
  \newcommand\tablename{Table}
\fi
}
\@ifpackageloaded{float}{}{\usepackage{float}}
\floatstyle{ruled}
\@ifundefined{c@chapter}{\newfloat{codelisting}{h}{lop}}{\newfloat{codelisting}{h}{lop}[chapter]}
\floatname{codelisting}{Listing}
\newcommand*\listoflistings{\listof{codelisting}{List of Listings}}
\makeatother
\makeatletter
\makeatother
\makeatletter
\@ifpackageloaded{caption}{}{\usepackage{caption}}
\@ifpackageloaded{subcaption}{}{\usepackage{subcaption}}
\makeatother
\makeatletter
\@ifpackageloaded{sidenotes}{}{\usepackage{sidenotes}}
\@ifpackageloaded{marginnote}{}{\usepackage{marginnote}}
\makeatother
\makeatletter
\@ifpackageloaded{tikz}{}{\usepackage{tikz}}
\makeatother
        \newcommand*\circled[1]{\tikz[baseline=(char.base)]{
          \node[shape=circle,draw,inner sep=1pt] (char) {{\scriptsize#1}};}}  
                  
\usepackage{bookmark}
\IfFileExists{xurl.sty}{\usepackage{xurl}}{} % add URL line breaks if available
\urlstyle{same}
% fallback for those not using the hyperref driver hyperxmp:
\makeatletter
\@ifundefined{xmpquote}{\newcommand{\xmpquote}[1]{#1}}{}
\makeatother
\hypersetup{
  pdftitle={260\_typst\_cheat.qmd},
  pdfauthor={JR},
  colorlinks=true,
  linkcolor={blue},
  filecolor={Maroon},
  citecolor={Blue},
  urlcolor={Blue},
  pdfcreator={LaTeX via pandoc}}


\title{260\_typst\_cheat.qmd}
\author{JR}
\date{2026-08-23}
\begin{document}
\maketitle


PURPOSE: Mixed quarto/typst features. Includes quarto div, html tricks
(legacy)

\subsection{TODO}\label{todo}

\begin{itemize}
\tightlist
\item
  create quarto template that embeds this setup
\item
  learn typst options for \_quarto.yml, under format: typst
\end{itemize}

\subsection{typst variable to hold
{[}content{]}}\label{typst-variable-to-hold-content}

\subsection{typst variables can hold {[}content{]}, or strings, lengths,
font sizes, numbers,
binary\ldots.}\label{typst-variables-can-hold-content-or-strings-lengths-font-sizes-numbers-binary.}

\subsection{PURPOSE:}\label{purpose}

This is \emph{.qmd file with (1) various quarto display tricks and (2)
}typst* chunks. /

\subsection{USE:}\label{use}

quarto render 260\_typst\_cheat.qmd --to typst

\subsection{TODO:}\label{todo-1}

Cleanup, remove remnants of past SCSS, remove what does \textbf{not}
work.

USE the \emph{typst app} ! - For typst \& math, use this file For Math.
- For typst \& Text, use \ldots{} .qmd - Typst for pdf (future html?) -
latex does seem to work. But plenty does NOT work on this page. - BEST:
use typst app

Watch \url{https://www.youtube.com/watch?v=2DbuqYKOsrY}\\
Tutorial: \url{https://typst.app/docs/tutorial/writing-in-typst/}

\begin{center}\rule{0.5\linewidth}{0.5pt}\end{center}

\begin{center}\rule{0.5\linewidth}{0.5pt}\end{center}

\begin{description}
\item[Def Word A]
Word A means \ldots..
\item[Word B]
B means this.
\end{description}

\begin{center}\rule{0.5\linewidth}{0.5pt}\end{center}

For CSS, use:

This is shadowbox in generic `div' container (Note: many ways to put
into .qmd) SEE:
https://developer.mozilla.org/en-US/docs/Web/HTML/Element/div

SEE: https://pandoc.org/chunkedhtml-demo/8.18-divs-and-spans.html

Should also be in shadobx, using pandox's colons

Latex, MATH, \textbf{diagrams}: SEE \texttt{code\_publish} esp MATH and
LATEX (but poor organization)

Embed latex.sty and real latex in quarto
https://youtu.be/2DbuqYKOsrY?t=396

Or, add \texttt{quarto\ extensions} (format: my\_extension) To add latex
or HTML forms \ldots{} Or, some R packages (gt, \ldots) allow you to
embed latex in a function. OR, typst (since 2022) - simpler than latex

\subsection{typst}\label{typst}

*.qmd, then template (position things like \(title\)), then style doc to
set title = ``Hello'' Typt is more English; compare to SCSS to make,
say, a box, yellow, with a font \ldots{}

\chapter{Section I}\label{section-i}

go to {[}Jim{]}

\section{TO DO}\label{to-do}

In SCSS file, there are - tags Example: h1 \{ \ldots{} \} - class
Example: .important \{ \ldots{} \} - id Example: \#title \{ \ldots\}

Inside the \{ \ldots{} \} is SCSS

Once defined, the SCSS can applied in *.qmd file in several ways

\begin{itemize}
\tightlist
\item
  chunk:
\end{itemize}

\begin{verbatim}
#| class-output:  myclass
\end{verbatim}

\begin{itemize}
\item
  text: add highlight \{.highlight\}
\item
  div:
\end{itemize}

\begin{verbatim}
::: {.border}
 text
:::
\end{verbatim}

\begin{itemize}
\item
  inline: {render this text}
\item
  via markdown: \emph{this is italics} this is underline
\end{itemize}

this is also underlined, using definition in SCSS

\begin{center}\rule{0.5\linewidth}{0.5pt}\end{center}

\section{To Publish quarto to
connect.posit.cloud}\label{to-publish-quarto-to-connect.posit.cloud}

https://connect.posit.cloud/jimrothstein

This is IDE (https://posit.cloud/content/yours?sort=name\_asc)

\begin{itemize}
\item
  Copy desired *.qmd to this folder: quarto \ldots{} test
\item
  Run the writeManifest() (below)
\item
  (Does render work locally?)
\item
  Check the resulting manifest.json is not bloated with excessive
  packages/dependencies.
\item
  commit/push
\item
  On this website: - https://connect.posit.cloud/jimrothstein
\item
  \texttt{Publish}
\item
  Choose Github, this folder, which *.qmd file
\end{itemize}

In this folder/project,

\begin{Shaded}
\begin{Highlighting}[]
\CommentTok{\#rsconnect::writeManifest()}
\end{Highlighting}
\end{Shaded}

\section{Absolute positioning
example}\label{absolute-positioning-example}

{I love .small}

{centering}

{[}absolute positioning{]}\{.absolute top=``50\%'' left-``50\%''\}

{and cake.}

\section{Preliminaries}\label{preliminaries}

\begin{tcolorbox}[enhanced jigsaw, arc=.35mm, bottomrule=.15mm, breakable, colback=white, colframe=quarto-callout-color-frame, left=2mm, leftrule=.75mm, opacityback=0, rightrule=.15mm, toprule=.15mm]

KEEP as QUARTO .qmd file (to learn to decorate) RStudio also - easy to
post completed quarto to Posit Cloud

\end{tcolorbox}

\subsection{REF}\label{ref}

\begin{itemize}
\item
  \url{https://sass-lang.com/documentation/}
\item
  \href{https://sass-lang.com/documentation/}{sass-lang.org}
\item
\item
\end{itemize}

Here is \href{https://sass-lang.com/documentation/}{FOO} and here is
\href{https://pandoc.org/MANUAL.html\#inline-links}{Pandoc links}

\begin{tcolorbox}[enhanced jigsaw, arc=.35mm, bottomrule=.15mm, bottomtitle=1mm, breakable, colback=white, colbacktitle=quarto-callout-tip-color!10!white, colframe=quarto-callout-tip-color-frame, coltitle=black, left=2mm, leftrule=.75mm, opacityback=0, opacitybacktitle=0.6, rightrule=.15mm, title=\textcolor{quarto-callout-tip-color}{\faLightbulb}\hspace{0.5em}{References, click to expand}, titlerule=0mm, toprule=.15mm, toptitle=1mm]

\begin{itemize}
\tightlist
\item
  Mimi - quarto, some CSS, SCSS (see her actual code):
  https://mine-cetinkaya-rundel.github.io/quarto-tip-a-day/\\
\item
  Themes, SASS:
\item
  SASS R pkg:
  https://cran.r-project.org/web/packages/sass/vignettes/sass.html\\
\item
  Quarto Themes + modify with SCSS:
  https://quarto.org/docs/output-formats/html-themes.html\#basic-options
  ~
\end{itemize}

\end{tcolorbox}

``Sass theme files can define both variables that globally set things
like colors and fonts, as well as rules that define more fine grained
behavior. To customize an existing Bootstrap theme with your own set of
variables and/or rules, just provide the base theme and then your custom
theme file(s):''

\begin{verbatim}
PURPOSE: Collect misc stuff related to quarto and HTML

- quarto 
- css, SCSS
- revealjs   (SEE:  examples in code_publish)
- HTML
\end{verbatim}

\subsection{Misc quarto tips}\label{misc-quarto-tips}

\begin{itemize}
\tightlist
\item
  Quarto looks to \_quarto.yml in current dir, then in parent dir.
\item
  This takes dominance over any scsc file in yaml in the .qmd file.
  (check !!)
\item
  \texttt{div} means container
\item
  \texttt{:::\ blocks} are called \emph{fences} and complies to
  \textless div \ldots\textgreater{}
\end{itemize}

https://www.sitepoint.com/understanding-and-using-rem-units-in-css/

\begin{itemize}
\tightlist
\item
  .em - height of entire metal type block, not the character alone.
\item
  .rem - size relative to the HTML root tag (\textasciitilde{} 16 px)
\end{itemize}

\begin{center}\rule{0.5\linewidth}{0.5pt}\end{center}

These are same: \texttt{::::\{.X\}} and \texttt{::::\{X\}} where X is
defined in CSS/SCSS.

\subsection{Output chunks}\label{output-chunks}

Note: use of \texttt{class-source} and \texttt{class-output}, define
.custombox\_* in SCSS class-source not working.

\begin{Shaded}
\begin{Highlighting}[]
\NormalTok{x}\OtherTok{=}\DecValTok{1}
\NormalTok{y }\OtherTok{=}\NormalTok{ x }\SpecialCharTok{+}\DecValTok{2}
\NormalTok{y}
\NormalTok{[}\DecValTok{1}\NormalTok{] }\DecValTok{3}
\end{Highlighting}
\end{Shaded}

\begin{Shaded}
\begin{Highlighting}[]
\NormalTok{x}\OtherTok{=}\DecValTok{1}
\NormalTok{y }\OtherTok{=}\NormalTok{ x }\SpecialCharTok{+}\DecValTok{2}
\NormalTok{y}
\NormalTok{[}\DecValTok{1}\NormalTok{] }\DecValTok{3}
\end{Highlighting}
\end{Shaded}

\begin{Shaded}
\begin{Highlighting}[]
\NormalTok{x}\OtherTok{=}\DecValTok{1}
\NormalTok{y }\OtherTok{=}\NormalTok{ x }\SpecialCharTok{+}\DecValTok{2}
\NormalTok{y}
\NormalTok{[}\DecValTok{1}\NormalTok{] }\DecValTok{3}
\end{Highlighting}
\end{Shaded}

\begin{Shaded}
\begin{Highlighting}[]
\CommentTok{\#| code{-}fold: true}
\CommentTok{\#| code{-}summmary: Show code.}

\NormalTok{Show me}
\end{Highlighting}
\end{Shaded}

\section{Pandoc only (cli) - pdf,}\label{pandoc-only-cli---pdf}

pandoc -D latex \# template pandoc -D html

templates look like \(var\), or \(IF\) etc. trick: Regardless of
content, place 3 back ticks at top/bottom will render document as
straight text.

\footnotesize

To render ascii document with small font (for printing):

\begin{itemize}
\item
  in , use ``` for verbatum
\item
  very top: \footenotesize
\item
  pandoc -f markdown -o out.pdf
\item
  pandoc -f markdown -V geometry:margin=2cm -o out.pdf
\item
  quarto render --to gfm --output
  \textless output\_file.md\textgreater{}
\item
  Visual Mode = menu driven, simplifies decorattion
\item
  Renders automatically.
\end{itemize}

\section{Quarto - basics}\label{quarto---basics}

markdown does NOT offer all tools (ex: change color) best to avoid
direct HTML insertion, use tools:

\begin{itemize}
\item
  knitr:: can be useful functions:
  \url{https://cran.r-project.org/web/packages/knitr/knitr.pdf}
\item
  rmarkdown cookbook:
  \url{https://bookdown.org/yihui/rmarkdown-cookbook}
\item
  markdown:: adds to rmarkdown:
\item
  Mimi: videos
\item
  Pandoc's version of markdown:
\item
  Bootstrap: predefined SEE getbootstrap.com css help-classes
\end{itemize}

\subsection{quarto examples}\label{quarto-examples}

SEE namespace file

\subsection{Run quarto (CLI) create project to create any quarto
project, including slides with
revealjs.}\label{run-quarto-cli-create-project-to-create-any-quarto-project-including-slides-with-revealjs.}

For revealjs, choose presentation when create document (in Rstudio)

\subsection{Quarto - markup, callouts}\label{quarto---markup-callouts}

\ul{underline text, where defined?}

{text with border, defined?}

\subsection{Create a footnote:}\label{create-a-footnote}

Need footnote one here: \footnote{This is footnote one.}

\begin{center}\rule{0.5\linewidth}{0.5pt}\end{center}

\section{Call out examples}\label{call-out-examples}

REF: https://quarto.org/docs/authoring/callouts.html

\subsection{.callout-note}\label{callout-note}

\begin{tcolorbox}[enhanced jigsaw, arc=.35mm, bottomrule=.15mm, bottomtitle=1mm, breakable, colback=white, colbacktitle=quarto-callout-note-color!10!white, colframe=quarto-callout-note-color-frame, coltitle=black, left=2mm, leftrule=.75mm, opacityback=0, opacitybacktitle=0.6, rightrule=.15mm, title=\textcolor{quarto-callout-note-color}{\faInfo}\hspace{0.5em}{Note}, titlerule=0mm, toprule=.15mm, toptitle=1mm]

This is .callout-note:\\
Note that there are five types of callouts, including: \texttt{note},
\texttt{warning}, \texttt{important}, \texttt{tip}, and
\texttt{caution}.

\end{tcolorbox}

\subsection{.callout-tip}\label{callout-tip}

\begin{tcolorbox}[enhanced jigsaw, arc=.35mm, bottomrule=.15mm, bottomtitle=1mm, breakable, colback=white, colbacktitle=quarto-callout-tip-color!10!white, colframe=quarto-callout-tip-color-frame, coltitle=black, left=2mm, leftrule=.75mm, opacityback=0, opacitybacktitle=0.6, rightrule=.15mm, title=\textcolor{quarto-callout-tip-color}{\faLightbulb}\hspace{0.5em}{Tip}, titlerule=0mm, toprule=.15mm, toptitle=1mm]

This is an example of a callout with a title.

\end{tcolorbox}

\subsection{.callout-caution, collapse}\label{callout-caution-collapse}

\begin{tcolorbox}[enhanced jigsaw, arc=.35mm, bottomrule=.15mm, bottomtitle=1mm, breakable, colback=white, colbacktitle=quarto-callout-caution-color!10!white, colframe=quarto-callout-caution-color-frame, coltitle=black, left=2mm, leftrule=.75mm, opacityback=0, opacitybacktitle=0.6, rightrule=.15mm, title=\textcolor{quarto-callout-caution-color}{\faFire}\hspace{0.5em}{Caution}, titlerule=0mm, toprule=.15mm, toptitle=1mm]

Expand To Learn About Collapse

This is an example of a `folded' caution callout that can be expanded by
the user. You can use \texttt{collapse="true"} to collapse it by default
or \texttt{collapse="false"} to make a collapsible callout that is
expanded by default.

\end{tcolorbox}

\subsection{.content-visible
when-format=html}\label{content-visible-when-formathtml}

\begin{tcolorbox}[enhanced jigsaw, arc=.35mm, bottomrule=.15mm, bottomtitle=1mm, breakable, colback=white, colbacktitle=quarto-callout-tip-color!10!white, colframe=quarto-callout-tip-color-frame, coltitle=black, left=2mm, leftrule=.75mm, opacityback=0, opacitybacktitle=0.6, rightrule=.15mm, title=\textcolor{quarto-callout-tip-color}{\faLightbulb}\hspace{0.5em}{Expand for Full Code Walkthrough}, titlerule=0mm, toprule=.15mm, toptitle=1mm]

\begin{Shaded}
\begin{Highlighting}[]
\FunctionTok{library}\NormalTok{(admiral)}
\FunctionTok{library}\NormalTok{(tibble)}
\NormalTok{adbds }\OtherTok{\textless{}{-}} \FunctionTok{tribble}\NormalTok{(}
  \SpecialCharTok{\textasciitilde{}}\NormalTok{PARAMCD, }\SpecialCharTok{\textasciitilde{}}\NormalTok{AVAL,}
  \StringTok{"AST"}\NormalTok{,    }\DecValTok{42}\NormalTok{,}
  \StringTok{"AST"}\NormalTok{,    }\DecValTok{52}\NormalTok{,}
  \StringTok{"AST"}\NormalTok{,    }\ConstantTok{NA\_real\_}\NormalTok{,}
  \StringTok{"ALT"}\NormalTok{,    }\DecValTok{33}\NormalTok{,}
  \StringTok{"ALT"}\NormalTok{,    }\DecValTok{51}
\NormalTok{)}
\CommentTok{\# Create a criterion flag with values "Y" and NA}
\FunctionTok{derive\_vars\_crit\_flag}\NormalTok{(}
\NormalTok{  adbds,}
  \AttributeTok{condition =}\NormalTok{ AVAL }\SpecialCharTok{\textgreater{}} \DecValTok{50}\NormalTok{,}
  \AttributeTok{description =} \StringTok{"Absolute value \textgreater{} 50"}
\NormalTok{)}
\CommentTok{\# Create criterion flag with values "Y", "N", and NA and parameter dependent}
\CommentTok{\# criterion description}
\FunctionTok{derive\_vars\_crit\_flag}\NormalTok{(}
\NormalTok{  adbds,}
  \AttributeTok{crit\_nr =} \DecValTok{2}\NormalTok{,}
  \AttributeTok{condition =}\NormalTok{ AVAL }\SpecialCharTok{\textgreater{}} \DecValTok{50}\NormalTok{,}
  \AttributeTok{description =} \FunctionTok{paste}\NormalTok{(PARAMCD, }\StringTok{"\textgreater{} 50"}\NormalTok{),}
  \AttributeTok{values\_yn =} \ConstantTok{TRUE}\NormalTok{,}
  \AttributeTok{create\_numeric\_flag =} \ConstantTok{TRUE}
\NormalTok{)}
\end{Highlighting}
\end{Shaded}

\end{tcolorbox}

\subsection{aside}\label{aside}

\marginnote{\begin{footnotesize}

\begin{enumerate}
\def\labelenumi{\arabic{enumi}.}
\tightlist
\item
  An aside.\\
\end{enumerate}

\end{footnotesize}}

\marginnote{\begin{footnotesize}

\begin{enumerate}
\def\labelenumi{\arabic{enumi}.}
\setcounter{enumi}{1}
\tightlist
\item
  Another aside;
\item
  contains a hidden keep a secret? USE \{.hidden\}
\end{enumerate}

\end{footnotesize}}

\subsection{hidden}\label{hidden}

4 This is secret

\begin{center}\rule{0.5\linewidth}{0.5pt}\end{center}

\subsection{Two columns}\label{two-columns}

I would like to have text here

Sentence becomes longer, it should automatically stay in their column

and here

More text

\begin{center}\rule{0.5\linewidth}{0.5pt}\end{center}

\begin{center}\rule{0.5\linewidth}{0.5pt}\end{center}

\subsection{Quarto \& CSS \textbar{} SCSS (see \_quarto.yml cosmo +
custom.scss)}\label{quarto-css-scss-see-_quarto.yml-cosmo-custom.scss}

in YAML, I used \texttt{date:\ \ today}\\
but in inline r 23 August 2026

1 verbatim:

\begin{verbatim}
in chunks: 
- OLD way: `{r style="xxx"; class=""} 
- NEW way: `{r .name } where .name is defined in .CSS or .SCSS 
\end{verbatim}

\begin{tcolorbox}[enhanced jigsaw, arc=.35mm, bottomrule=.15mm, bottomtitle=1mm, breakable, colback=white, colbacktitle=quarto-callout-note-color!10!white, colframe=quarto-callout-note-color-frame, coltitle=black, left=2mm, leftrule=.75mm, opacityback=0, opacitybacktitle=0.6, rightrule=.15mm, title=\textcolor{quarto-callout-note-color}{\faInfo}\hspace{0.5em}{Note}, titlerule=0mm, toprule=.15mm, toptitle=1mm]

2nd way:

in chunks: - OLD way: `\{r style=``xxx''; class=``\,``\}

\begin{itemize}
\tightlist
\item
  NEW way: `\{r .name \} where .name is defined in .CSS or .SCSS
\end{itemize}

\end{tcolorbox}

\begin{center}\rule{0.5\linewidth}{0.5pt}\end{center}

Defined in SCSS (.custombox\_white)

\begin{itemize}
\tightlist
\item
  Div Examples, defined in prior CSS block as \texttt{greybox}
\end{itemize}

\begin{center}\rule{0.5\linewidth}{0.5pt}\end{center}

\subsection{Custom Blocks
(rmarkdown-9.6)}\label{custom-blocks-rmarkdown-9.6}

\marginnote{\begin{footnotesize}

(aside) Custom Blocks

\end{footnotesize}}

\protect\phantomsection\label{hello}
Hello \textbf{world}! is this CSS or SCSS?

\begin{center}\rule{0.5\linewidth}{0.5pt}\end{center}

\begin{verbatim}
This example ... renders as:
:::{.small}
Is this .small?
:::
\end{verbatim}

Is this .small?

\begin{center}\rule{0.5\linewidth}{0.5pt}\end{center}

\begin{itemize}
\item
  BE SURE to use \texttt{\_quarto.yml} in base directory.
\item
  \texttt{\_quarto.yml} identifies css or SCSS file.
\item
  Quarto uses \texttt{themes} from Bootstrap, in special file format (*
  \ldots. *) OR, use existing SCSS
\item
  Because \texttt{\textless{}}, \texttt{\textgreater{}} have meaning in
  HTML, use \texttt{escaped} version \texttt{\textless{}\&} in HTTP
\end{itemize}

\begin{verbatim}


::: aside 
SEE References 
:::

Interactive? javascript widgets or **Shiny** Fiddle with code? use
**WebR**

Decorate?\
- Use assorted tags that compile to html div - In chunks, use knitr to
decorate that chunk

------------------------------------------------------------------------


### SCSS examples (mostly)

***

``` {.small}
1. This is small b/c I used `.small` 
\end{verbatim}

\begin{enumerate}
\def\labelenumi{\arabic{enumi}.}
\setcounter{enumi}{1}
\tightlist
\item
  But this time, \texttt{small} is inside the colon blocks
\end{enumerate}

\begin{Shaded}
\begin{Highlighting}[]
\NormalTok{3. Very small}
\end{Highlighting}
\end{Shaded}

\begin{enumerate}
\def\labelenumi{\arabic{enumi}.}
\setcounter{enumi}{3}
\tightlist
\item
  Inside colon blocks
\end{enumerate}

\begin{center}\rule{0.5\linewidth}{0.5pt}\end{center}

\begin{Shaded}
\begin{Highlighting}[]



\FunctionTok{\#\#\# Display code and chunk options}


\NormalTok{::: \{.cell\}}

\InformationTok{\textasciigrave{}\textasciigrave{}\textasciigrave{}\textasciigrave{}\{.cell{-}code  code{-}line{-}numbers="|5|6"\}}
\NormalTok{\textasciigrave{}\textasciigrave{}\textasciigrave{}}\OperatorTok{\{}\NormalTok{r}\OperatorTok{\}}
\PreprocessorTok{\#}\ErrorTok{| label: example}
\PreprocessorTok{\#}\ErrorTok{| code{-}line{-}numbers: "|5|6"}
\PreprocessorTok{\#}\ErrorTok{| echo: fenced \# shows these options in output}
\NormalTok{x  }\OperatorTok{\textless{}{-}} \DecValTok{1}
\NormalTok{y  }\OperatorTok{\textless{}{-}} \DecValTok{2}
\NormalTok{x}\OperatorTok{+}\NormalTok{y}
\end{Highlighting}
\end{Shaded}

````

\begin{verbatim}
[1] 3
\end{verbatim}

:::

\begin{center}\rule{0.5\linewidth}{0.5pt}\end{center}

TODO pdf - fix greater than equa pdf: toc: false include-in-header:
text: \textbar{}

\usepackage[top=.5in, bottom=.5in, left=.7in, right=.5in]{geometry}
        \usepackage{amsmath,amssymb,amstext,amsfonts}
        \usepackage{lipsum}

Will only appear in PDF. \renewcommand{\familydefault}{\sfdefault}

\section{R Definitions \& Examples   (TERSE)}
\section{Manuals}

\textbackslash section(Definitions\}

\begin{center}\rule{0.5\linewidth}{0.5pt}\end{center}

\subsection{table with scss}\label{table-with-scss}

First Name

Last Name

Super Hero

Peter

Parker

Spiderman

Bruce

Wayne

Batman

Clark

Kent

Superman

\begin{center}\rule{0.5\linewidth}{0.5pt}\end{center}

\begin{center}\rule{0.5\linewidth}{0.5pt}\end{center}

\begin{center}\rule{0.5\linewidth}{0.5pt}\end{center}

Example of Bootstrap classes:

\begin{Shaded}
\begin{Highlighting}[]
\NormalTok{mtcars[}\DecValTok{1}\SpecialCharTok{:}\DecValTok{5}\NormalTok{, }\StringTok{"mpg"}\NormalTok{]}
\NormalTok{[}\DecValTok{1}\NormalTok{] }\FloatTok{21.0} \FloatTok{21.0} \FloatTok{22.8} \FloatTok{21.4} \FloatTok{18.7}
\end{Highlighting}
\end{Shaded}

To make sure that we always get a data frame, we have to use the
argument \texttt{drop\ =\ FALSE}. Now we use the chunk option
\texttt{class.source\ =\ "bg-success"}.

\begin{Shaded}
\begin{Highlighting}[]
\NormalTok{mtcars[}\DecValTok{1}\SpecialCharTok{:}\DecValTok{5}\NormalTok{, }\StringTok{"mpg"}\NormalTok{, drop }\OtherTok{=} \ConstantTok{FALSE}\NormalTok{]}
\NormalTok{                   mpg}
\NormalTok{Mazda RX4         }\FloatTok{21.0}
\NormalTok{Mazda RX4 Wag     }\FloatTok{21.0}
\NormalTok{Datsun }\DecValTok{710}        \FloatTok{22.8}
\NormalTok{Hornet }\DecValTok{4}\NormalTok{ Drive    }\FloatTok{21.4}
\NormalTok{Hornet Sportabout }\FloatTok{18.7}
\end{Highlighting}
\end{Shaded}

\begin{center}\rule{0.5\linewidth}{0.5pt}\end{center}

\subsection{Annotate code and create code to change font
color:}\label{annotate-code-and-create-code-to-change-font-color}

Use knitr to create for red font, rmarkdown cookbook 5.1

\protect\phantomsection\label{annotated-cell-17}%
\begin{Shaded}
\begin{Highlighting}[]
\NormalTok{col }\OtherTok{=} \ControlFlowTok{function}\NormalTok{(text, color)\{ }\hspace*{\fill}\NormalTok{\circled{1}}
    \ControlFlowTok{if}\NormalTok{ (knitr}\SpecialCharTok{::}\FunctionTok{is\_html\_output}\NormalTok{()) \{ }\hspace*{\fill}\NormalTok{\circled{2}}
        \FunctionTok{sprintf}\NormalTok{(}\StringTok{"\textless{}span style=\textquotesingle{}color:\%s\textquotesingle{}\textgreater{} \%s \textless{}/span\textgreater{}"}\NormalTok{, color, text )}
\NormalTok{    \}}
\NormalTok{\}}
\end{Highlighting}
\end{Shaded}

\begin{description}
\tightlist
\item[\circled{1}]
Annotate line 1
\item[\circled{2}]
Annotate line 2
\end{description}

\begin{center}\rule{0.5\linewidth}{0.5pt}\end{center}

\subsection{Fancy Theorems}\label{fancy-theorems}

\begin{verbatim}
not working

::: column-margin
We know from *the first fundamental theorem of calculus* that for $x$ in
$[a, b]$:

$$\frac{d}{dx}\left( \int_{a}^{x} f(u)\,du\right)=f(x).$$

:::
\end{verbatim}

\section{R chunks with decorations}\label{r-chunks-with-decorations}

\subsection{Box (rmarkdown-book 9.6)}\label{box-rmarkdown-book-9.6}

\begin{Shaded}
\begin{Highlighting}[]
\NormalTok{mtcars[}\DecValTok{1}\SpecialCharTok{:}\DecValTok{3}\NormalTok{, }\DecValTok{1}\SpecialCharTok{:}\DecValTok{2}\NormalTok{]}
\NormalTok{               mpg cyl}
\NormalTok{Mazda RX4     }\FloatTok{21.0}   \DecValTok{6}
\NormalTok{Mazda RX4 Wag }\FloatTok{21.0}   \DecValTok{6}
\NormalTok{Datsun }\DecValTok{710}    \FloatTok{22.8}   \DecValTok{4}
\end{Highlighting}
\end{Shaded}

\begin{Shaded}
\begin{Highlighting}[]
\NormalTok{mtcars[}\DecValTok{1}\SpecialCharTok{:}\DecValTok{3}\NormalTok{, }\DecValTok{1}\SpecialCharTok{:}\DecValTok{2}\NormalTok{]}
\NormalTok{               mpg cyl}
\NormalTok{Mazda RX4     }\FloatTok{21.0}   \DecValTok{6}
\NormalTok{Mazda RX4 Wag }\FloatTok{21.0}   \DecValTok{6}
\NormalTok{Datsun }\DecValTok{710}    \FloatTok{22.8}   \DecValTok{4}
\end{Highlighting}
\end{Shaded}

\begin{center}\rule{0.5\linewidth}{0.5pt}\end{center}

\protect\phantomsection\label{special}
This does what?

\section{CSS here}\label{css-here}

\subsection{CSS in-line examples.}\label{css-in-line-examples.}

PURPOSE: Collect of HTML and PDF inline commands to control output

\subsection{in-file, CSS examples}\label{in-file-css-examples}

\subsection{Quarto: In-line CSS Examples for HTML,
PDF}\label{quarto-in-line-css-examples-for-html-pdf}

Very small

\subsection{Inline}\label{inline}

We can apply styles to a sentence or a word by creating spans using
\texttt{{[}{]}} to surround the sentence or word that we want to style
and use \texttt{\{\}} to define the style that we want to apply. For
example,

The color of this word is {red}. And {this line has a yellow
background}.

\begin{center}\rule{0.5\linewidth}{0.5pt}\end{center}

\subsection{Create a CSS class
(inline)}\label{create-a-css-class-inline}

\begin{verbatim}
create class "myColor" (red) inline | use in paragraph | and div block| 
\end{verbatim}

Example 1

myColor

Example 2

This is a Div defined by CSS, inline

\subsubsection{Header uses in-file CSS
class}\label{header-uses-in-file-css-class}

\begin{Shaded}
\begin{Highlighting}[]
\FunctionTok{head}\NormalTok{(mtcars, }\AttributeTok{n=}\DecValTok{2}\NormalTok{)}
\NormalTok{              mpg cyl disp  hp drat    wt  qsec vs am gear carb}
\NormalTok{Mazda RX4      }\DecValTok{21}   \DecValTok{6}  \DecValTok{160} \DecValTok{110}  \FloatTok{3.9} \FloatTok{2.620} \FloatTok{16.46}  \DecValTok{0}  \DecValTok{1}    \DecValTok{4}    \DecValTok{4}
\NormalTok{Mazda RX4 Wag  }\DecValTok{21}   \DecValTok{6}  \DecValTok{160} \DecValTok{110}  \FloatTok{3.9} \FloatTok{2.875} \FloatTok{17.02}  \DecValTok{0}  \DecValTok{1}    \DecValTok{4}    \DecValTok{4}
\end{Highlighting}
\end{Shaded}

\section{Acknowledgments}\label{acknowledgments}

I am grateful for the insightful comments offered by the anonymous peer
reviewers at Books \& Texts. The generosity and expertise of one and all
have improved this study in innumerable ways and saved me from many
errors; those that inevitably remain are entirely my own responsibility.




\end{document}
